\section{Conclusion}\label{sec:conclusion}

\rt{} replaces the summed group statistics of the sums-based \clamp{} design by one secret-shared scalar per client, sorted, trimmed and privatised by two non-colluding servers. We proved that up to $b$ falsifying clients cannot move the released estimate outside the range of honest gaps, that the release is attribute-level differentially private with a sensitivity fixed by public validation-set sizes, and a high-probability bound on the released gap of the returned model. In experiments with four fixed attacks and an adaptive one, up to $b$ liars did not move the test gap away from its honest level, where one liar had disabled the sums-based \clamp{}. With ten releases that honest level was close to training without control on the Pokec graphs; with per-round releases the controller was active, lowered the gap modestly, and liars with $f\le b$ again did not change the outcome. Liars that also poisoned their updates added little to the update attack when $f\le b$ and broke the gap channel when $f>b$. A full release round cost 1.1--10.5\,s in a semi-honest two-party implementation, dominated by the norm checks on the updates.

The design has clear costs and limits. With noise on a single trimmed scalar and ten releases, the honest-case gap on Credit and the Pokec graphs was higher than with the sums-based \clamp{}; on graphs with a few hundred nodes the noise exceeds the signal; robustness ends at $f=b+1$, where a hiding attack turns the controller against its purpose; fairness poisoning through updates is not countered; and the training experiments simulate the secure computation with continuous noise. Making the robust channel as accurate as the pooled one, and pairing it with a robust update channel, remain open.
